\section{Реализация}
\label{sec:Chapter6} \index{Chapter6}

В данной главе описывается практическая реализация SLP-векторизатора в составе компилятора Go. Рассматриваются среда разработки, ключевые структуры данных, алгоритмы каждого этапа, генерация векторного кода и методика тестирования.

\subsection{Инструментарий и среда разработки}
\label{sec:dev_env}

Разработка велась во внутреннем форке языка Go (согласованном с версией go1.22) с целевой архитектурой ARM64 с поддержкой SIMD-расширения NEON. Проход SLP реализован в пакете \texttt{cmd/compile/internal/ssa} в виде отдельного файла \texttt{slp.go}. Новый проход зарегистрирован в общем списке оптимизаций компилятора.

Для отладки использовалась встроенная в компилятор система флагов: переменная \texttt{debug}, управляемая флагом \texttt{-d=ssa/slp/debug=N}, а также переменная окружения \texttt{GOSSAFUNC}, позволяющая вывести SSA-граф конкретной функции до и после прохода. Это позволило визуально контролировать корректность преобразований.

Тестирование выполнялось с помощью системы \texttt{go test}, включая специализированные тесты \texttt{codegen}: \texttt{asmcheck} (проверка наличия/отсутствия ожидаемых векторных инструкций в ассемблерном листинге) и \texttt{errorcheck} (проверка диагностических сообщений для случаев, когда векторизация не должна применяться).

\subsection{Базовые структуры данных и проверки}
\label{sec:data_structures}

Основными структурами данных являются типы \texttt{pack} и \texttt{packSet}.

\begin{verbatim}
type pack []*Value
\end{verbatim}

\texttt{pack} определён как упорядоченный срез указателей на значения SSA. Все значения в одном \texttt{pack}'е изоморфны и упорядочены согласно линиям.

\begin{verbatim}
type packSet []pack
\end{verbatim}

\texttt{packSet} --- срез \texttt{pack}'ов, накапливаемый в процессе анализа блока.

Ключевые вспомогательные проверки, используемые на этапе анализа:

\begin{itemize}
    \item \texttt{isIsomorphic(v1, v2)} --- проверяет, что две инструкции имеют одинаковый код операции и одинаковый тип результата. Дополнительно требует равенства количества аргументов.
    
    \item \texttt{isIndependent(v1, v2)} --- проверяет, что инструкции не зависят друг от друга по данным и разделяют одно состояние памяти.
    
    \item \texttt{isAdjacent(v1, v2)} --- проверяет смежность двух обращений к памяти. Использует анализ из прохода \texttt{memcombine} для разложения указателя на базовую часть, индекс и смещение. Два обращения считаются смежными и упорядоченными, если у них совпадают база и индекс, а разность константных смещений равна размеру элемента.
    
    \item \texttt{canPack(v1, v2, Ps, b)} --- собирает все проверки для возможности образования новой пары: изоморфизм, векторизуемость операции и типа, смежность (для памяти), отсутствие конфликтов с существующими группами (ни одно значение уже не задействовано) и независимость.
\end{itemize}

\subsection{Проход \texttt{slp}: главный управляющий цикл}
\label{sec:main_loop}

Входная точка прохода --- функция \texttt{slp(f *Func)}. Она обходит все базовые блоки функции; для каждого блока \(b\) последовательно выполняются следующие шаги:

\begin{enumerate}
    \item \textbf{Инициализация} (\texttt{findSeeds}): формируется начальный набор из пар смежных загрузок/сохранений.
    \item \textbf{Расширение} (\texttt{extendPackSet}): итеративно вызываются \texttt{followUse} и \texttt{followDef} для добавления зависимых инструкций.
    \item \textbf{Комбинирование} (\texttt{combinePacks}): соседние группы сливаются.
    \item \textbf{Разбиение} (\texttt{splitPackSets}): набор групп разбивается на связные компоненты.
    \item Для каждой компоненты:
        \begin{itemize}
            \item валидация (\texttt{checkPackSet}) --- проверка корректности;
            \item планирование (\texttt{schedulePacks}) --- топологическая сортировка;
            \item легализация (\texttt{legalizePackSet}) --- приведение к аппаратной ширине 16 байт;
            \item генерация кода (\texttt{emitPacks}) --- создание векторных инструкций и замена скаляров.
        \end{itemize}
\end{enumerate}

При наличии нескольких компонент каждая обрабатывается независимо, что позволяет векторизовать несколько не связанных участков в одном блоке.

\subsection{Этапы анализа}
\label{sec:analysis_stages}

\subsubsection{Поиск начальных пар}

Процедура двойным вложенным циклом перебирает все значения блока. Для каждой пары \((v_1, v_2)\), где обе инструкции являются обращениями к памяти (\texttt{OpLoad} или \texttt{OpStore}), вызывается \texttt{canPack}. При успехе создаётся новая группа. Сложность этого шага квадратична от числа инструкций в блоке, но на практике размеры блоков невелики, и с учётом ранних отсечений (проверка кода операции и типа до вызова \texttt{canPack}) накладные расходы остаются приемлемыми.

\subsubsection{Расширение групп}

Выполняется итеративный цикл до сходимости: за каждую итерацию для всех текущих групп вызываются \texttt{followUse} и \texttt{followDef}.

\begin{itemize}
    \item \texttt{followDef} для группы \(\texttt{pack}\{v_1, v_2\}\) перебирает все аргументы; для каждого проверяет \texttt{canPack(v1.Args[i], v2.Args[i])} и добавляет новые группы при успехе. Добавление происходит для всех совместимых пар аргументов.
    \item \texttt{followUse} работает аналогично: он находит все использования \(v_1\) и \(v_2\) внутри текущего блока и проверяет, можно ли их сгруппировать.
\end{itemize}

\subsubsection{Комбинирование}

После расширения набор групп содержит только пары. Процедура ищет группы, такие что конец одной совпадает с началом другой, и объединяет их. Исходные мелкие группы заменяются объединённой. Процесс повторяется, пока возможны слияния.

\subsubsection{Разбиение на компоненты связности}

Алгоритм выделяет связные компоненты графа, вершинами которого являются группы, а ребро между ними существует, если какой-либо аргумент значения из одной группы принадлежит другой группе, или наоборот. Это даёт набор независимых компонент.

\subsection{Валидация и планирование}
\label{sec:validation2}

\subsubsection{Валидация}

Выполняет три основные проверки:

\begin{enumerate}
    \item \textbf{checkExternalUses} --- для каждого значения в группах ищутся все его использования в функции (\texttt{findAllUses}). Использование в том же наборе, а также особые случаи (\texttt{Store}-память для продолжения цепочки памяти) считаются допустимыми. Также для длинных цепочек разрешено ограниченное количество внешних использований. Любое другое внешнее использование приводит к отказу от векторизации всей компоненты.
    
    \item \textbf{checkClosedOperands} --- проверяет, что все аргументы принадлежат тому же набору групп. Нарушение замкнутости свидетельствует о зависимости, выходящей за пределы векторизуемого подграфа.
    
    \item \textbf{checkLaneCoherence} --- удостоверяется, что все группы в компоненте имеют одинаковую длину и что значения используются в тех же линиях, в которых они производятся. Это проверка когерентности линий; в противном случае потребовались бы \texttt{shuffle}-инструкции.
\end{enumerate}

\subsubsection{Планирование}

Реализована простая топологическая сортировка. Для каждой группы рассматриваются аргументы её первого элемента (поскольку все элементы изоморфны, структура аргументов одинакова). Пропуская указатели и память, алгоритм находит группу-аргумент. Если такая группа ещё не отмечена запланированной, текущая группа откладывается. Итерации продолжаются, пока все группы не будут запланированы.

\subsection{Генерация векторного SSA-кода}
\label{sec:codegen2}

\subsubsection{Легализация}

Перед генерацией кода группы разбиваются до ширины в 128 бит. Если ширина групп не кратна 128 бит, то весь набор отбрасывается.

\subsubsection{Создание векторных инструкций}

Для каждой группы создаётся ровно одна векторная инструкция. Для отображения используется глобальная хеш-таблица соответствия. Примеры:

\begin{itemize}
    \item \texttt{OpLoad} → \texttt{OpLoad} с типом результата \texttt{TypeVec128}.
    \item \texttt{OpStore} → \texttt{OpStore} с типом результата \texttt{TypeVec128}.
    \item \texttt{OpSub64} → \texttt{OpSubInt64x2} (векторное вычитание).
    \item \texttt{OpCopy} → \texttt{OpCopy} с типом \texttt{TypeVec128}.
\end{itemize}

Для операций обращений к памяти адрес нагрузки берётся из первого элемента группы (левая граница), память --- с помощью метода \texttt{getMemoryArg}, возвращающего общий \texttt{Memory} для всей группы. Для операций, использующих другие группы, нужные аргументы находятся в кэше уже созданных векторных операций.

Все вновь созданные векторные значения сохраняются в \texttt{packToSimd} (отображение \texttt{pack} → \texttt{Value}), который служит кэшем для последующих групп.

\subsubsection{Замена скалярных использований}

После создания всех векторных инструкций итеративно обходятся исходные группы и для каждого значения обновляются его использования:

\begin{itemize}
    \item Если использование принадлежит другой группе (т.е. внутреннее), оно уже было подставлено на этапе формирования аргументов векторной инструкции.
    \item Для значений \texttt{Store} обновляется аргумент памяти у следующих за ними инструкций: вместо результата скалярного \texttt{Store} подставляется память от векторного \texttt{Store}.
    \item Для внешних использований вставляется \texttt{OpGetElem\{Type\}x\{Number\ of\ lanes\}} --- инструкция извлечения элемента из векторного регистра по индексу, соответствующему позиции значения в группе.
\end{itemize}

После замены исходные скалярные инструкции теряют всех пользователей и удаляются следующим проходом \texttt{deadcode}.

\subsection{Тестирование корректности}
\label{sec:testing}

Тестирование проводилось на двух уровнях:

\begin{itemize}
    \item \textbf{errorcheck-тесты} --- файлы с директивами \texttt{// errorcheck}, проверяющие, что SLP-проход не применяется в известных сложных случаях. В теле функций размещены аннотации вида \texttt{// ERROR "текст диагностики"}. При запуске компилятора с флагом \texttt{-d=ssa/slp/debug=1} векторизатор выводит сообщения с причиной отказа, и тестовая система сопоставляет их с ожидаемыми. Таким образом проверяются сценарии: разные состояния памяти, вызов между инструкциями, внешние использования, нечётная ширина и др.
    
    \item \textbf{asmcheck-тесты} --- проверяют наличие ожидаемых векторных инструкций в ассемблерном выводе для ARM64. В комментариях к функциям указаны регулярные выражения (например, \texttt{// arm64:\textbackslash VSUB\textbackslash t}), которые должны присутствовать в сгенерированном коде.
\end{itemize}

Тесты обеспечивают регрессионную защиту при дальнейшем развитии прохода и документируют его текущие возможности и ограничения. Ручное тестирование с помощью \texttt{GOSSAFUNC} позволяло визуально инспектировать SSA-граф до и после прохода, отслеживая корректность подстановок и отсутствие оборванных зависимостей.